%%%%%%%%%%%%%%%%%%%%%%%%%%%%%%%%%%%%%%%%%%%%%%%%%%%%%%%%%%%%%%%%%%%%%%%
%                                                                     %
%                        Chapter 2 - Section 1                        %
%                                                                     % 
%%%%%%%%%%%%%%%%%%%%%%%%%%%%%%%%%%%%%%%%%%%%%%%%%%%%%%%%%%%%%%%%%%%%%%%


\chapter{Successive Approximations}
\index{successive approximation}

	
In this chapter we continue exploring the mathematical implications of
the $S$-$I$-$R$ model.  In the last chapter we calculated future
values of $S$, $I$, and $R$ by assuming that the rates $S'$, $I'$, and
$R'$ stayed fixed for a whole day.  Since the rates are {\em not\/}
fixed---they change with $S$, $I$, and $R$---the values of $S$, $I$,
and $R$ we obtained have to be considered as estimates only.  In this
chapter we will see how to build a succession of better and better
estimates that get us as close as we wish to the true values implied
by the model.

This method of {\bf successive approximation} is a basic tool of
calculus.  It is the one fundamentally new process you will encounter,
the ingredient that sets calculus apart from the mathematics you have
already studied.  With it you will be able to solve a vast array of
problems that other methods can't handle.


\section{Making Approximations}

In chapter 1 we looked at the specific $S$-$I$-$R$ model:
\begin{align*}
S' &= -.00001 \, SI, \\
I' &= .00001 \, SI - I/14, \\
R' &= I/14,
\end{align*}
with initial values at time $t = 0$:  
\[
S = 45400,  \quad I = 2100, \quad R = 2500.
\]

We originally developed this model as a {\em description} of the
relations among the different components of an epidemic.  Almost
immediately, though, we began using the rate equations in the model as
a {\em recipe} for predicting what happens over the course of the
epidemic: If we know at some time $t$ the values of $S(t), I(t),
\mbox{ and } R(t)$), \mar{Rate equations tell us where to go next}
then the equations tell us how to estimate values of the functions at
other times. We used this approach in the last chapter to move
backwards and forwards in time, calculating the values of $S$, $I$,
and $R$ as we went.

While we got numbers, there were some questions about how accurate
these numbers were---that is, how exactly they represented the values
implied by the model.  In the process we called ``there and back
again'' we used current values of $S$, $I$, and $R$ to find the rates,
used these rates to go forward one day, recalculate the rates, and
come back to the present---and we got different values from the ones
we started with!  Resolving this discrepancy will be an important
feature of the technique developed in this section.

\subsection{The Longest March Begins with a Single Step}

So far, in generating numbers from the $S$-$I$-$R$ rate equations, we
have assumed that the rates remained constant over an entire day, or
longer. Since the rates aren't constant---they depend on the values of
$S$, $I$, and $R$, which are always changing---the values we
calculated for the variables at times other than the given initial
time are, at best, estimates.  These estimates, while incorrect, are
not useless. Let's see how they behave in the ``there and back again''
process of chapter 1 as we recalculate the rates more and more
frequently, producing a sequence of approximations to the values we
are looking for.

\subsubsection{There and Back Again Again}
\index{there and back again}

On page~\pageref{`thereback'} in chapter 1 we used the rate equations to
go forward a day and come back again.  We started with the initial
values
\[
S(0)=45400, \quad I(0) = 2100, \quad R(0) = 2500,
\]
calculated the rates, went forward a day to $t=1$, recalculated the
rates, and came back a day to $t=0$. We ended up with the estimates
\[
S(0)=45737.1, \quad I(0) = 1820.3, \quad R(0) = 2442.6
\]
---which are rather far from the values of $S(0)$, $I(0)$, and $R(0)$
we started with.

A clue to the resolution of this discrepancy appeared in problem 18,
page~\pageref{`therebackprob'}.  There you were asked to go forward two
days and come back again in two different ways, using $\Delta t=2$ (a
total of 2 steps) in the first case and $\Delta t=1$ (a total of 4
steps) in the second.  Here are the resulting values calculated for
$S(0)$ in each case and the discrepancy between this value and the
original value $S(0)=45400$:
\begin{center}
\begin{tabular}{ccc}
step size & new $S(0)$ & discrepancy \\[2pt] \hline
$\Delta t=2$ & 46717.6 & 1317.6\\
$\Delta t=1$ & 46021.3 & 621.3 \\
\end{tabular}
\end{center}

While the discrepancy is fairly large in either case, $\Delta t=1$
clearly does better than $\Delta t=2$.  But if smaller
%
\mar{Smaller steps generate a smaller discrepancy}
%
is better, why stop at $\Delta t=1$?  What happens if we take even
smaller time steps, get the corresponding new values of $S$, $I$, and
$R$, and use these values to recalculate the rates each time?

Recall that the rate $S'$ (or $I'$ or $R'$) is simply the multiplier
which gives $\Delta S$---the (estimated) change in $S$---for a given
change $\Delta t$ in $t$
\[
\Delta S = S'\cdot \Delta t\, .
\]
This relation holds for any value of $\Delta t$, integer or not.  Once
we have this value for $\Delta S$, we can then calculate
\[
\mbox{new (estimated) } S = \mbox{current (estimated) } S + \Delta S
\]
in the usual way.  Note that we have written ``(estimated)''
throughout to emphasize the fact that if $S'$ is not constant over the
entire time $\Delta t$, then the value we get for $\Delta S$ will
typically be only an approximation to the real change in $S$.

Let's try going forward one day and coming back, using different
values for $\Delta t$.  As we reduce $\Delta t$ the number of
calculations will increase.  The program SIR we used in the last
chapter can still be used to do the tedious calculations.  Thus if we
decide to use 10 steps of size $\Delta t=.1$, we would just change two
lines in that program: 
\vspace*{4pt}\\
\hspace*{.5in}{\tt deltat = .1 \\
\hspace*{.5in}FOR k = 1 TO 10 \vspace*{4pt}\\} 
If we now run SIR with these modifications we can verify the following
sequence of values (The values have been rounded off, and the {\tt
  PRINT} statement has been modified to show the new values of the
rates at each step as well):
\hspace*{-.25in}
\begin{center}
	Estimated values of $S, I, \mbox{ and } R$ \\ [-1pt]
	for step sizes $\Delta t = .1$
\medskip
\small

\hspace*{-.25in}\begin{tabular}{ccccccc}
	$t$ &  $S(t)$  & $I(t)$ & $R(t)$ & $S'(t)$ & $I'(t)$ & $R'(t)$\\ 
   [2pt]\hline
		\begin{tabular}{l}
		0.0  \rule{0pt}{14pt}   \\
		0.1  \\
		0.2		\\
		0.3	\\
		\multicolumn{1}{c}{\vdots}  \\
		1.0
		\end{tabular}
	& 
		\begin{tabular}{r}
		\rule{0pt}{14pt} 45\,400.0  \\
		45\,304.7	\\
		45\,205.9	\\
		45\,103.6  \\
		\multicolumn{1}{c}{\vdots}	\\
		44\,278.7	
		\end{tabular} 
	&
		\begin{tabular}{r}
		\rule{0pt}{14pt} 2100.0 \\
		2180.3 \\
		2263.6  \\
		2349.7  \\
		\multicolumn{1}{c}{\vdots}  \\
		3042.9
		\end{tabular}
	&
		\begin{tabular}{r}
		\rule{0pt}{14pt} 2500.0 \\
		2515.0 \\
		2530.6  \\
		2546.7  \\
		\multicolumn{1}{c}{\vdots}  \\
		2678.4		
		\end{tabular}
	&
		\begin{tabular}{r}
		\rule{0pt}{14pt} --953.4 \\
		--987.8  \\
		--1023.3  \\
		--1059.8  \\
		\multicolumn{1}{c}{\vdots} \\
		--1347.7
		\end{tabular}
	&
		\begin{tabular}{r}
		\rule{0pt}{14pt} 803.4  \\
		832.1  \\
		861.6  \\
		892.0  \\
		\multicolumn{1}{c}{\vdots}	\\
		1130.0	
		\end{tabular}
	&
		\begin{tabular}{r}
		\rule{0pt}{14pt} 150.0  \\
		155.7  \\
		161.7  \\
		167.8  \\
		\multicolumn{1}{c}{\vdots}  \\
		217.4
		\end{tabular}
	\end{tabular}
\end{center}
\normalsize


Having arrived at $t=1$, we can now use SIR to turn around and go back
to $t=0$.  Here's how:
\begin{itemize}

\item Change the initial line of the program to {\tt t = 1} to reflect
  our new starting time.

\item Change the next three lines to use the values we just calculated
  for $S(1)$, $I(1)$, and $R(1)$ as our {\em starting} values in
  \mar{The {\em sign} of {\tt deltat} determines whether we move
    forward or backward in time} SIR.

\item Change the value of {\tt deltat} to be {\tt -.1} (so each time
  the program executes the command {\tt t = t + deltat} it reduces the
  value of {\tt t} by .1).

\end{itemize}
With these changes SIR will yield the desired estimates for $S(0)$,
$I(0)$, and $R(0)$, and we get
\[
S(0)=45433.5, \quad I(0) = 2072.3, \quad R(0) = 2494.3.
\]
This is clearly a considerable improvement over the values obtained
with $\Delta t=1$.

With this promising result, the obvious thing to do is to try even
smaller values of $\Delta t$, perhaps $\Delta t = .01$.  We could
continue using SIR, making the needed modifications each
time. Instead, though, let's rewrite SIR \index{program!SIR} slightly
to make it better suited to our current needs.  Look at the program
SIRVALUE\index{program!SIRVALUE} below and compare it with SIR.

\newpage
   
\vspace*{4pt}
\noindent
\footnotesize
\begin{minipage}[t]{3in}
\begin{center}
{\bf Program: SIRVALUE}\label{sirvalue}
\end{center}
\vspace{-18pt}

\begin{verbatim}
tinitial = 0
tfinal = 1
t = tinitial
S = 45400
I = 2100
R = 2500
numberofsteps = 10
deltat = (tfinal - tinitial)/numberofsteps
FOR k = 1 TO numberofsteps
      Sprime = -.00001 * S * I
      Iprime = .00001 * S * I - I / 14
      Rprime = I / 14
      deltaS = Sprime * deltat
      deltaI = Iprime * deltat
      deltaR = Rprime * deltat
      t = t + deltat
      S = S + deltaS
      I = I + deltaI
      R = R + deltaR
NEXT k
PRINT t, S, I, R
\end{verbatim}
\end{minipage} \ \
\begin{minipage}[t]{2.8in}
\begin{center}
{\bf Program: SIR}
\end{center}
\vspace{-18pt}

\begin{verbatim}
t = 0
S = 45400
I = 2100
R = 2500
deltat = .1
FOR k = 1 TO 10
      Sprime = -.00001 * S * I
      Iprime = .00001 * S * I - I / 14
      Rprime = I / 14
      deltaS = Sprime * deltat
      deltaI = Iprime * deltat
      deltaR = Rprime * deltat
      t = t + deltat
      S = S + deltaS
      I = I + deltaI
      R = R + deltaR
      PRINT t, S, I, R
NEXT k
\end{verbatim}
\end{minipage}
\normalsize
\vspace{4pt}

You will see that the major change is to place the {\tt PRINT}
statement outside the loop, so only the final values of $S$, $I$, and
$R$ get printed.  This speeds up the work, since otherwise, with
$\Delta t = .001$, for instance, we would be asking the computer to
print out 1000 lines---about 30 screens of text!  Another change is
that the value of {\tt deltat} no longer needs to be specified---it is
automatically determined by the values of {\tt tinitial}, {\tt
  tfinal}, and {\tt numberofsteps}.

As written above, SIR and SIRVALUE both run for 10 steps of size
$0.1\, .$ By changing the value of {\tt numberofsteps} in the program
we can quickly get estimates
%	
\mar{With a computer we can generate lots of data and look for patterns}
%
for $S(1)$, $I(1)$, and $R(1)$ for a wide range of values for $\Delta
t \, .$ Moreover, once we have these estimates we can use SIRVALUE
again to go backwards in time to $t=0$, by making changes similar to
those we made in SIR earlier.  First, we need to change the value of
{\tt tinitial} to 1 and the value of {\tt tfinal} to 0.  Notice that
this automatically will make {\tt deltat} a negative quantity, so that
each time we run through the loop we step back in time.  Second, we
need to set the starting values of {\tt S, I,} and {\tt R} to the
values we just obtained for $S(1)$, $I(1)$, and $R(1)$.  With these
changes, SIRVALUE will give us the corresponding estimated values for
$S(0)$, $I(0)$, and $R(0)$.

If we use SIRVALUE with $\Delta t$ ranging from 1 to .00001 (which
means letting {\tt numberofsteps} range from 1 to 100,000) we get the
table below.  This table lists the computed values of $S(1)$, $I(1)$,
and $R(1)$ for each $\Delta t$, followed by the estimated value of
$S(0)$ obtained by running SIRVALUE backward in time from these new
values, and, finally, the discrepancy between this estimated value of
$S(0)$ and the original value $S(0) = 45400.$

\begin{center}
	Estimated values of $S, I, \mbox{ and } R$ when $t = 1$, \\ [1pt]
	for step sizes $\Delta t = 10^{-N}$, $N = 0, \ldots 5$, \\ [3pt]
	together with the corresponding backwards estimate for $S(0)$.\label{backwards2}
\bigskip
\footnotesize

	\hspace*{-.5in}\begin{tabular}{c|ccc|c|c}
	$\Delta t$ &  $S(1)$  & $I(1)$ & $R(1)$ & new $S(0)$ & discrepancy\\ 
   [2pt]\hline
		\begin{tabular}{l}
		1.0 \rule{0pt}{14pt} \\
		0.1		\\
		0.01		\\
		0.001	\\
		0.000\,1  \\
		0.000\,01	
		\end{tabular}
	& 
		\begin{tabular}{l}
		44\,446.6 \rule{0pt}{14pt} \\
		44\,278.6648	\\
		44\,257.8301	\\
		44\,255.6960  \\
		44\,255.4821  \\
		44\,255.4607
		\end{tabular} 
	&
		\begin{tabular}{l}
		2903.4 \rule{0pt}{14pt} \\
		3042.9241 \\
		3060.1948  \\
		3061.9633  \\
		3062.1406  \\
		3062.1584
		\end{tabular}
	&
		\begin{tabular}{l}
		2650.0 \rule{0pt}{14pt} \\
		2678.4111 \\
		2681.9751  \\
		2682.3406  \\
		2682.3773 \\
		2682.3809
		\end{tabular}
	&
		\begin{tabular}{l}
		45\,737.0626 \rule{0pt}{14pt} \\
		45\,433.4741  \\
		45\,403.3615  \\
		45\,400.3363  \\
		45\,400.0336  \\
		45\,400.0034
		\end{tabular}
	&
		\begin{tabular}{r}
		\rule{0pt}{14pt} 337.0626  \\
		33.4741  \\
		3.3615  \\
		.3363  \\
		.0336  \\
		.0034
		\end{tabular}
	\end{tabular}
\end{center}
\normalsize
\medskip
There are several striking features of this table.  The first is that
if we go forward one day and come back again, we can get back as close
as 
%
\mar{Smaller steps generate a discrepancy which can be made as
  small as we like} 
%
we want to our initial value of $S(0)$ {\em provided we recalculate
  the rates frequently enough}.  After 200,000 rounds of calculations
($\Delta t = .00001$) we ended up only .0034 away from our starting
value.  In fact, there is a clear pattern to the values of the errors
as we decrease the step size.  In the exercises it is left for you to
explore this pattern and show that similar results hold for $I$ and
for $R$.

A second feature is that as we read down the column under $S(1)$, we
find each digit {\bf stabilizes}\index{stabilize}---that is, after
changing for a while, it eventually becomes fixed at a particular
value.  The initial digits 44 are the first to stabilize, and that
happens by the time $\Delta t = 0.1$.  Then the third digit 2
stabilizes, when $\Delta t = 0.01$.  Roughly speaking, one more digit
stabilizes at each successive level.  The table is revealing to us,
digit by digit, the true value of $S(1)$.  By the fifth stage we learn
that the integer part of $S(1)$ is 44255.  By the sixth stage we can
say that the true value of $S(1)$ is $44255.4\ldots$ .

When we write $S(1) = 44255.4\ldots$ 
%
\mar{Approximations lead to exact values} 
%
we are expressing $S(1)$ to {\bf one decimal place
  accuracy}.\index{accuracy to a given number of decimal places} This
says, first, that the decimal expansion of $S(1)$ begins with exactly
the six digits shown and, second, that there are further digits after
the 4 (represented by the three dots ``\ldots'')\label{`dotdot'}.  In
this case, we can identify further digits simply by continuing the
table.  Since our step sizes have the form $\Delta t = 10^{-N}$, we
just need to increase $N$.  For example, to express $S(1)$ accurately
to six decimal places, we need to stabilize the first eleven digits in
our estimates of $S(1)$.  The table suggests that $\Delta t$ should
probably be about $10^{-10}$---i.e., $N = 10$.

The true value of $S(1)$ emerges through a process that generates a
sequence of successive approximations.\index{successive approximation}
We say $S(1) = 44255.4$\ldots is the {\bf limit}\index{limit} of this
sequence as $\Delta t$ is made smaller and smaller or, equivalently,
as $N$ is made larger and larger.  We also say that the sequence of
successive approximations {\bf converges}\index{convergence!of a
  sequence} to the limit $S(1)$.  Here is a mathematical notation that
expresses these statements more compactly:
\begin{align*}
	S(1) &= \lim_{\Delta t \to 0}
   \{\mbox{the estimate of $S(1)$}\} \qquad \mbox{or, equivalently,}\\
	&= \lim_{N \to \infty}
		\{\mbox{the estimate of $S(1)$ 
		with $\Delta t = 10^{-N}$}\}.
\end{align*}
The symbol $\infty$ stands for ``infinity,'' and the expression $N \to
\infty$ is often pronounced ``as $N$ goes to infinity.''  However, it
is often more instructive to say ``as $N$ gets larger and larger,
without bound.''

You should check that similar patterns are occurring in the $I(1)$ and
$R(1)$ columns as well.

\aside{The limit concept lies at the heart of calculus.  Later on
  we'll give a precise definition, but you should first see limits at
  work in a number of contexts and begin to develop some intuitions
  about what they are.  This approach mirrors the historical
  development of calculus---mathematicians freely used limits for well
  over a century before a careful, rigorous definition was developed.}

\subsection{One Picture Is Worth a Hundred Tables}

As we noted, the program SIRVALUE prints out only the final values of
$S$, $I$, and $R$ because it would typically take too much space to
print out the intermediate values.  However, if instead of printing
these values we plot them graphically, we can convey all this
intermediate information in a compact and comprehensible form.

Suppose, for instance, that we wanted to record the
calculations leading up to $S(3)$ by plotting all the points.  The
graphs below plot all the pairs of values $(t, S)$ that are
calculated along the way for the cases $\Delta t =1$ (4 points),
$\Delta t =.1$ (31 points), and $\Delta t =.01$ (301
points).

\newpage

\mbox{}

\vspace{240pt}

\special{psfile=../ps/calc1/picture3.ps}

By the time we get to steps of size .01, the resulting graph
begins to look like a continuous curve.  This suggests that instead of
simply plotting the points we might want to draw lines connecting the
points as they're calculated.


We can easily modify SIRVALUE to do this---the only
changes will be to replace the {\tt PRINT} command with a command to
draw a line and to move this command inside the loop (so that it is
executed every time new values are computed).  We will also need to
add a line or two at the beginning to tell the computer to set up the
screen to plot points.  This usually involves opening a
\index{window}{\bf window}---i.e., specifying the horizontal and
vertical ranges the screen should depict. Since programming languages
vary slightly in the way this is done, we use italicized text ``{\sl
  Set up GRAPHICS\/}'' to make clear that this statement is {\bf not}
part of the program---you will have to express this in the form your
programming language specifies.  Similarly, the command
\begin{quote} 
\rule{6pt}{0pt} {\sl Plot the line from {\tt (t, S)}
      to {\tt (t + deltat, S + deltaS)}}
\end{quote}
will have
to be stated in the correct format for your language. The
computational core of SIRVALUE is unchanged.  Here is what the new
program looks like if we want to use $\Delta t = .1$ and connect the
points with straight lines:

\newpage

\begin{center}
{\bf Program: SIRPLOT}\index{program!SIRPLOT}\label{sirplot}
\end{center}
{\sl Set up GRAPHICS}
\vspace{-10pt}
\begin{verbatim}
   tinitial = 0
   tfinal = 3
   t = tinitial
   S = 45400
   I = 2100
   R = 2500
   numberofsteps = 30
   deltat = (tfinal - tinitial)/numberofsteps
   FOR k = 1 TO numberofsteps
         Sprime = -.00001 * S * I
         Iprime = .00001 * S * I - I / 14
         Rprime = I / 14
         deltaS = Sprime * deltat
         deltaI = Iprime * deltat
         deltaR = Rprime * deltat
\end{verbatim}

\vspace{-8pt}

\rule{32pt}{0pt} {\sl Plot the line from 
		{\tt (t, S)} to {\tt (t + deltat, S + deltaS)}}

\vspace{-12pt}

\begin{verbatim}
         t = t + deltat
         S = S + deltaS
         I = I + deltaI
         R = R + deltaR
   NEXT k
\end{verbatim}

If we had wanted just to plot the points, we could have used a command
of the form {\sl Plot the point} {\tt (t, S)} in place of the command
to plot the line, moving this command down two lines so it came after
we had computed the new values of {\tt t} and {\tt S}.  We would also
need to place that command before the loop so that the initial point
corresponding to $t = 0$ gets plotted.

When we ``connect the dots'' like this we emphasize graphically the
underlying assumption we have been making in all our estimates: that
the function $S(t)$ is linear (i.e., it is changing at a constant
rate) over each interval $\Delta t$.  Let's see what the graphs look
like when we do this for the three values of $\Delta t$ we used above.
To compare the results more readily we'll plot the graphs on the same
set of axes. (We will look at a program for doing this in the next
section.) We get the following picture:

\newpage

\mbox{}

\vspace*{200pt}

\special{psfile=../ps/calc1/picture2.ps}

The graphs become indistinguishable from each other and increasingly
 look like smooth
%
\mar{Graphs made up of\\ line segments look like smooth curves if\\ the
 segments are\\ short enough}
%
curves as the number of segments increases.  If we plotted the
 3000-step graph as well, it would be indistinguishable from the
 300-step graph at this scale.  If we now shift our focus from the end
 value $S(3)$ and look at all the intermediate values as well, we find
 that each graph gives an approximate value for $S(t)$ for {\em every}
 value of $t$ between 0 and 3.  We are seeing the entire function
 $S(t)$ over this interval.

Just as we wrote
\[
S(3) = \lim_{N \to \infty}\{\mbox{the estimate of $S(3)$ 
		with $\Delta t = 10^{-N}$}\}.
\]
We can also write 
\[
\mbox{graph of } S(t) =\lim_{N \to \infty}
		\{\mbox{line-segment approximations with $\Delta t = 
		10^{-N}$}\}.
\] 
The way we see the graph of $S(t)$ emerging from successive
 approximations is our first example of a fundamental result. It has
 wide-ranging implications which will occupy much of our attention for
 the rest of the course.

\newpage

\subsection{Piecewise Linear Functions}
\index{piecewise linear function}

Let's examine the implications of this approach more closely by
considering the ``one-step" ($\Delta t=3$) approximation to $S(t)$ and
the ``three-step" ($\Delta t=1$) approximation over the time interval
$0\leq t\leq 3$. In the first case we are making the simplifying
assumption that $S$ decreases at the rate $S' = -953.4$ persons per
day for the entire three days.  In the second case we use three
shorter steps of length $\Delta t=1$, with the slopes of the
corresponding segments given by the table on page~\pageref{`threeday'}
in Chapter~1, summarized below (note that since $\Delta t = 1$ day we
have that the magnitude of $\Delta S = S'\cdot\Delta t$ is the same as
the magnitude of $S'$):
\begin{center}
\begin{tabular}{ccc}
	$t$ & $S$ &  $S'$\\ [1pt] \hline
	\begin{tabular}{r} \rule{0pt}{14pt} 0 \\ 1 \\ 2 \\ 3 
		\end{tabular} &
	\begin{tabular}{r} 
		\rule{0pt}{14pt} $45\,400.0$ \\ $44\,446.6$ \\
 		$43\,156.1$ \\ $41\,435.7$ \end{tabular} &
	\begin{tabular}{r} 
		\rule{0pt}{14pt} $-953.4$ \\ $-1\,290.5$ \\ 
		$-1\,720.4$ \\ \mbox{} \end{tabular} 
\end{tabular}
\end{center}
Here are the corresponding graphs we get:
\begin{center}
\begin{picture}(340,232)(-55,-32)
	\put(0,0){\vector(1,0){250}}
	\put(50,0){\line(0,1){3}}
	\put(100,0){\line(0,1){3}}
	\put(150,0){\line(0,1){3}}
	\put(0,-12){\makebox[0pt]{0}}
	\put(50,-12){\makebox[0pt]{1}}
	\put(100,-12){\makebox[0pt]{2}}
	\put(150,-12){\makebox[0pt]{3}}
	\put(245,-12){\makebox[0pt]{days}}
	\put(245,5){\makebox[0pt]{$t$}}
	\put(0,0){\line(0,1){15}}
	\put(0,22){\circle*{2}}
	\put(0,30){\circle*{2}}
	\put(0,38){\circle*{2}}
	\put(0,45){\vector(0,1){155}}
	\put(-4,191){\makebox[0pt][r]{persons}}
	\put(5,191){\makebox[0pt][l]{$S$}}
	\put(-4,174){\makebox[0pt][r]{45400}}
	\begin{thicklines}
		\put(0,175){\line(2,-1){150}}
		\put(50,150){\line(3,-2){50}}
		\put(100,116.7){\line(6,-5){50}}
	\end{thicklines}
	\put(0,175){\circle*{4}}
	\put(50,150){\circle*{4}}
	\put(150,100){\circle*{4}}
	\put(100,116.7){\circle*{4}}
	\put(150,75){\circle*{4}}
	\multiput(0,175)(5,0){31}{\circle*{1}}
	\multiput(0,150)(5,0){10}{\circle*{1}}
	\multiput(50,116.7)(5,0){10}{\circle*{1}}
	\multiput(100,75)(5,0){10}{\circle*{1}}
	\multiput(150,100)(0,5){15}{\circle*{1}}
	\multiput(50,116.7)(0,5){7}{\circle*{1}}
	\multiput(100,75)(0,5){8}{\circle*{1}}
	\put(-15,162){\vector(1,0){13}}
	\put(-18,158){\makebox[0pt][r]{$-953.4$}}
	\put(48,129){\makebox[0pt][r]{$-1290.5$}}
	\put(98,90){\makebox[0pt][r]{$-1720.4$}}
	\put(75,178){\makebox[0pt]{3 days}}
	\put(156,133){\makebox[0pt][l]{$-953.4 \times 3$}}
	\put(156,96){\makebox[0pt][l]{42539.8}}
	\put(156,71){\makebox[0pt][l]{41435.7}}
	\put(125,182){\makebox[0pt][l]{\bf one large step}}
	\put(78,58){\makebox[0pt][l]{\bf three small steps}}
	\put(122,-32){\makebox[0pt]
		{Two approximations to $S$ during the first three 
		days}}
\end{picture}
\end{center}

\textbf{The ``one-step'' estimate}.  Assuming that $S$ decreases at
 the rate $S' = -953.4$ persons per day for the entire three days is
 equivalent to assuming that $S$ follows the upper graph---a straight
 line with slope \mbox{--953.4} persons/day.  In other words, the
 one-step approach approximates $S$ by a {\bf linear
 function}\index{linear function} of $t$.  If we use the notation
 $S_{1}(t)$ to denote this (one-step) linear approximation we
 have \begin{quote} One-step estimate: \qquad $S(t) \approx S_{1}(t)=
 45400 - 953.4 \, t \, .$
	\end{quote}
Because the one-step estimate is actually a function we can find the
 value of $S_{1}(t)$ for {\em all\/} $t$ in the interval $0 \leq t
 \leq 3$, not just $t=3$, and thereby get corresponding estimates for
 $S(t)$ as well.  For example, 
\begin{align*}
S_{1}(2) &= 45400 - 953.4 \times 2 = 43493.2 \\ 
S_{1}(1.7) &= 45400 - 953.4 \times 1.7 = 43779.22 \, .
\end{align*}

\textbf{The ``three-step'' estimate}.  With three smaller steps of
 size $\Delta t = 1$, we get a function whose graph is composed of
 three line segments, each starting at the $(t, S)$ point at the
%
\mar{A single function\\ may not be specified\\ by a single equation}
%
beginning of each day and with a slope equal to the corresponding rate
 of new infections at the beginning of the day.  Let's call this
 function $S_{3}(t)$. The three-step estimate $S_{3}(t)$ is hence not
 a linear function, strictly speaking.  However, since its graph is
 made up of several straight pieces, it is called a {\bf piecewise
 linear function}.\index{piecewise linear function} Recalling that the
 equation of a line through the point $(x_0, y_0)$ with slope $m$ can
 be written in the form $y=m(x-x_0)+y_0$, we can use the values for
 $S$ (which correspond to the $y$-values) and $S'$ (which give us the
 slopes of the segments) at times $t=0$, $t=1$, and $t=2$ calculated
 above to get an explicit formula for $S_{3}(t)$:
\[
S_3(t)=\left\{\begin{array}{ll}
			y = -953.4(t-0)+45400 & \mbox{ if $0\leq t\leq 1$} \\
			y=-1290.5(t-1)+44446.6&\mbox{ if $1\leq t\leq 2$} \\
			y=-1720.4(t-2)+43156.1&\mbox{ if $2\leq t\leq 3$}
		\end{array}
	\right.
\]

Note that we have used $\leq$ in the defining formulas since at the
 values $t=1$ and $t=2$ it doesn't matter which equation we use.  This
 is equivalent to saying that the straight line segments are connected
 to each other.  Since the slopes of the three segments of $S_{3}(t)$
 are progressively more negative, the piecewise linear graph gets
 progressively steeper as $t$ increases.  This explains why the value
 $S_{3}(3)$ is {\em lower\/} than the value of $S_{1}(3)$.  While it
 is rare that we would actually need to write down an explicit formula
 like this for the piecewise-linear approximation---it is easier, and
 usually more informative, just to define $S_{3}(t)$ by its graph---it
 is nevertheless important to realize that there really is an
 approximating function defined for all values of $t$ in the interval
 $[0, 3]$, not just at the finite set of $t$ values where we make the
 recalculations.
	
By the time we are dealing with the 300-step function $S_{300}(t)$ we
 can't even tell by its graph that it is piecewise linear unless we
 zoom in very close.  In principle, though, we could still write down
 a simple linear formula for each of its segments (see the exercises).

\medskip

\textbf{An appraisal}.  The graph of $S_{1}(t)$ gives us a rough idea
 of what is happening to the true function $S(t)$ during the first
 three days.  It starts off at the same rate as $S'$, but subsequently
 the rates move apart.  The value of $S_{1}'$ never changes, while
 $S'$ changes with the (ever-changing) values of $S$ and~$I$.

The graph of $S_{3}(t)$ is a distinct improvement because it changes
 its direction twice, modifying its slope at the beginning of each day
 to come back into agreement with the rate equation.  But since the
 three-step graph is still piecewise linear, it continues to suffer
 from the same shortcoming as the one-step: once we restrict our
 attention to a single straight segment (for example, where $1 \leq t
 \leq 2$), then the three-step graph {\em also\/} has a constant
 slope, while $S'$ is always changing.  Nevertheless, $S_{3}(t)$ does
 satisfy the rate equation in our original model three times---at the
 beginning of each segment---and isn't too far off at other
 times. When we get to $S_{300}(t)$ we have a function which satisfies
 the rate equation at 300 times and is very close in between.
	
Each of these graphs gives us an idea of the behavior of the true
 function $S(t)$ during the time interval $0 \leq t \leq 3$.  None is
 strictly correct, but none is hopelessly wrong, either.  All are {\bf
 approximations} to the truth.  Moreover, $S_{3}(t)$ is a {\bf better}
 approximation than $S_{1}(t)$---because it reflects at least some of
 the variability in $S'$---and $S_{300}(t)$ is better still.  Thus,
 even before we have a clear picture of the shape of the true function
 $S(t)$, we would expect it to be closer to $S_{300}(t)$ than to
 $S_{3}(t)$. As we saw above, when we take piecewise linear
 approximations with smaller and smaller step sizes, it is reasonable
 to think that they will approach the true function $S$ in the limit.
Expressing this in the notation we have used before,
\[
\mbox{the function } S(t) =\lim_{N \to \infty} 
   \{\mbox{the chain of linear functions with $\Delta t = 10^{-N}$}\}.
\]

\newpage
	
\subsection{Approximate versus Exact}

You may find it unsettling that our efforts give us only a sequence of
 approximations to $S(3)$, and not the exact value, or only a sequence
 of piecewise-linear approximations to $S(t)$, not the ``real''
 function itself.  In what sense can we say we ``know'' the number
 $S(3)$ or the function $S(t)$?  The answer is: in
%
\mar{What does it mean\\ to "know" a number like~$\pi$?}
%
the same sense that we ``know'' a number like $\sqrt{2}$ or $\pi$.
  There are two distinct aspects to the way we know a number.  On the
 one hand, we can {\bf characterize} a number precisely and
 completely:

\smallskip
\begin{tabular}{rl}
	$\pi$: & the ratio of the circumference of a circle to its diameter; \\
	$\sqrt{2}$: & the positive number whose square is 2; \\
\end{tabular}

\smallskip
\noindent  
On the other hand, when we try to {\bf construct} the decimal
 expansion of a number, we usually get only approximate and incomplete
 results.  For example, when we do calculations by hand we might use
 the rough estimates $\sqrt{2} \approx 1.414$ and $\pi \approx
 3.1416$.  With a desk-top computer we might have $\sqrt{2} \approx
 1.414\,213\,562\,373\,095$ and $\pi \approx
 3.141\,592\,653\,589\,793$, but these are still approximations, and
 we are really saying
\[
\begin{array}{rcl}
	\sqrt{2} &=& 1.414\,213\,562\,373\,09 \ldots \\
	\pi &=& 3.141\,592\,653\,589\,79 \ldots \ .
	\end{array}
\]
The {\em complete\/} decimal expansions for $\sqrt{2}$ and $\pi$ are
 unknown!  The exact values exist as limits of approximations that
 involve successively longer strings of digits, but we never see the
 limits---only approximations. In the final section of this chapter we
 will see ways of generating these approximations for $\sqrt{2}$ and
 for $\pi$.

What we say about $\pi$ and $\sqrt{2}$ is true for $S(3)$ in exactly
 the same way.  We can {\em characterize\/} it quite precisely, and we
 can {\em construct\/} approximations to its numerical value to any
 desired degree of accuracy.  Here, for example, is a characterization
 of $S(3)$:
\begin{quote}
The $S$-$I$-$R$ problem for which $a = .00001$ and $b = 1/14$ and for
 which $S = 45400$, $I = 2100$, $R = 2500$ when $t = 0$ determines
 three functions $S(t)$, $I(t)$, and $R(t)$.  The number $S(3)$ is the
 value that the function $S(t)$ has when $t = 3$.
\end{quote} 
		
You should try to extend this argument to describe the sense in which
 we ``know'' the function $S(t)$ by knowing its piecewise-linear
 approximations. Try to convince yourself that this is operationally
 no different from the way we ``know'' functions like $f(x)=\sqrt{x}$.
  In each instance we can {\bf characterize} the function completely,
 but we can only {\bf construct} an approximation to most values of
 the function or to its graph.
	
All this discussion of approximations may strike you as an unfortunate
 departure from the accuracy and precision you may have been led to
 expect
%
\mar{Being able to approximate a number to 12 decimal places is
 usually as good as knowing its value precisely}
%
in mathematics up until now.  In fact, it is precisely this ability to
 make quick and accurate approximations to problems that is one of the
 most powerful features of mathematics.  This is what goes on every
 time you use your calculator to evaluate $\log 3$ or $\sin 37$.  Your
 calculator doesn't really know what these numbers are---but it does
 know how to approximate them quickly to 12 decimal places.  Similar
 kinds of approximations are also at the heart of how bridges are
 built and spaceships are sent to the moon.

{\bf A Caution:} The fact that computers and calculators are really
 only dealing with approximations
%
\mar{However \ldots \\ \mbox{}}
%
when we think they are being exact occasionally leads to problems, the
 most common of which involves {\bf roundoff errors}.\index{roundoff
 error} You can probably generate a relatively harmless manifestation
 of this on your computer with the SIRVALUE program.  Modify the {\tt
 PRINT} line so it prints out the final value of {\tt t} to 10 or 12
 digits, and try running it with a high value for {\tt numberofsteps},
 say 1 million or 10 million.  You would expect the final value of
 {\tt t} to be exactly 1 in every case, since you are adding {\tt
 deltat = 1/numberofsteps} to itself {\tt numberofsteps} times.  The
 catch is that the computer doesn't store the exact value {\tt
 1/numberofsteps} unless {\tt numberofsteps} is a power of 2.  In all
 other cases it will only be using an approximation, and if you add up
 enough quantities that are slightly off, their cumulative error will
 begin to show.  We will encounter a somewhat less benign
 manifestation of roundoff error in the next chapter.


\subsection{Exercises}

\noindent
\textbf{There and back again}
\index{there and back again}

\vspace{-5pt}

\numa Look at the table on page~\pageref{backwards2}.  What is your
 best guess of the {\em exact} value of $I(1)$?  (Use the ``\ldots ''
 notation introduced on page~\pageref{`dotdot'}.)

\sub  What is the exact value of $R(1)$?

\num We noted that the discrepancy (the difference between the new
 estimate for $S(0)$ and the original value) seemed to decrease as
 $\Delta t$ decreased.  \sub What is your best estimate (using only
 the information in the table) for the value of $\Delta t$ needed to
 produce a discrepancy of .001 ?  \sub More generally, express as
 precisely as you can the apparent relation between the size of the
 discrepancy and the size of $\Delta t$.

\numa Suppose you wanted to try going three days forward and then
 coming back, using $\Delta t=.01$.  What changes would you have to
 make in SIRVALUE to do this?  

\sub Make a table similar to the one on page~\pageref{backwards2} for
 going three days forward and coming back for $\Delta t=1$, .1, .01,
 and.001.

\sub In this new table how does the size of the discrepancy for a
 given value of $\Delta t$ compare with the value in the original
 table?

\sub What value of $\Delta t$ do you think you would need to determine
 the integer parts of $S(3)$ and $R(3)$ exactly?

\medskip

\textbf{Piecewise linear functions}
\index{piecewise linear function}

\vspace{-5pt} 

\num Using this three-step approximation, what is $S_{3}(1.7)$?  What
 is $S_{3}(2.5)$?  \hfill

\num How would you modify SIRVALUE to get $S_{3000}(3)$ ?  Do it; what
 do you get?

\num What additional changes would you make to get the values of $t$,
 $S$, and $S'$ at the beginning of the 193rd segment of $S_{300}(t)$ ?
  [HINT: You only need to alter the {\tt FOR k = 1 TO numberofsteps}
 line (since you don't want to go all the way to the end) and the {\tt
 PRINT} line.  (Note that after running the loop for, say, 20 times,
 the values of {\tt t} and {\tt S} are the values for the beginning of
 the 21st segment, while the value of {\tt Sprime} will still be the
 slope of the 20th segment.)]

\num  Suppose we wanted to determine the value of $S_{300}(2.84135)$.

\sub In which of the 300 segments of the graph of $S_{300}(t)$ would
 we look to find this information?

\sub What are the values of $t$, $S$, and $S'$ at the beginning of
 this segment?

\sub  What is the equation of this segment?

\sub What is $S_{300}(2.84135)$ ?

\num How would you modify SIRVALUE to calculate estimates for $S$,
 $I$, and $R$ when $t = -6$, using $\Delta t=.05$ ?  Do it; what do
 you get?


\num We want to use SIRPLOT to look at the graph of $S(t)$ over the
 first 20 days, using $\Delta t=.01$.

\sub  What changes would we have to make in the program?

\sub  Sketch the graph you get when you make these changes.

\sub If you wanted to plot the graph of $I(t)$ over this same time
 interval, what additional modifications to SIRPLOT would be needed?
  Make them, and sketch the result.  When does the infection appear to
 hit its peak?

\sub Modify SIRPLOT to sketch on the same graph all three functions
 over the first 70 days.  Sketch the result.

\medskip

\textbf{The DO--WHILE loop}
\index{loop!DO-WHILE}
\label{DO-WHILE}

A difficulty in giving a precise answer to the last question was that
 we had to get all the values for 20 days, then go back to estimate by
 eye when the peak occurred.  It would be helpful if we could write a
 program that ran until it reached the point we were looking for, and
 then stopped.  To do this, we need a different kind of loop---a {\em
 conditional\/} loop that keeps looping only while some specified
 condition is true.  A DO--WHILE loop is one useful way to do
 this. Here's how the modified SIRPLOT program would look:
\vspace*{.1in}
\index{program!DO-WHILE}

\enlargethispage*{20pt}
{\small
{\sl Set up GRAPHICS}
\vspace{-4pt}
\begin{verbatim}
   tinitial = 0
   t = tinitial
   S = 45400
   I = 2100
   R = 2500
   Iprime = .00001 * S * I - I / 14
   deltat = .01
   DO WHILE Iprime > 0
         Sprime = -.00001 * S * I
         Iprime = .00001 * S * I - I / 14
         Rprime = I / 14
         deltaS = Sprime * deltat
         deltaI = Iprime * deltat
         deltaR = Rprime * deltat
\end{verbatim}
\vspace{-4pt}
\rule{47pt}{0pt} {\sl Plot the line from 
		{\tt (t, S)}} {\sl to 
		{\tt (t + deltat, S + deltaS)}}\\
\vspace{-20pt}
\begin{verbatim}
         t = t + deltat
         S = S + deltaS
         I = I + deltaI
         R = R + deltaR
   LOOP
   PRINT t - deltat
\end{verbatim}
}%end \small

\noindent
The changes we have made are:
\begin{itemize}

\item Since we don't know what the final time will be, eliminate the
  {\tt tfinal = 20} and the {\tt numberofsteps = 2000} commands.

\item Since the condition in our loop is keyed to the value of $I'$,
  we have to calculate the initial value of $I'$ before the loop
  starts.

\item Instead of {\tt deltat = (tfinal - tinitial)/numberofsteps} use
  the statement {\tt deltat = .01}.

\item The key change is to replace the {\tt FOR k = 1 TO
  numberofsteps} line by the line {\tt DO WHILE Iprime > 0}.

\item To denote the end of the loop we replace the {\tt NEXT k}
  command with the command {\tt LOOP}.

\item After the {\tt LOOP} command add the line {\tt PRINT t -
  deltat}.  (The reason we had to back up one step at the end is
  because due to the way the program was written, the computer takes
  one final step with a negative value for {\tt Iprime} before it
  stops.)

\end{itemize}

The net effect of all this is that the program will continue working
as before, calculating values and plotting points {\tt (t, S)}, but
{\em only as long as the condition in the {\tt DO WHILE} statement is
  true.}  The condition we used here was that $I'$ had to be
positive---this is the condition that ensures that values of $I$ are
still getting bigger.  While this condition is true, we can always get
a larger value for $I$ by going forward another increment $\Delta t$.
As soon as the condition is false---i.e., as soon as $I'$ is
negative---the values for $I$ will be decreasing, which means we have
passed the peak and so want to stop.

\num  Make these modifications; what value for $t$ do you get?

\num You could modify the {\tt PRINT t} command to also print out
other quantities.

\sub What is the value of $I$ at its peak?

\sub What is the value of $S$ when $I$ is at its peak?  Does this
agree with the threshold value we predicted in chapter 1?

\num Suppose you change the initial value of $S$ to be 5400 and run
the previous program.  Now what happens?  Why?

\num Suppose we wanted to know how long the epidemic lasts.  We could
use DO--WHILE and keep stepping forward using, say $\Delta t = .1$, so
long as $I \geq 1$.  As soon as $I$ was less than 1 we would want to
stop and see what the value of $t$ was.

\sub What modifications would you make in SIRVALUE to get this
information?

\sub Run your modified program.  What value do you get for $t$?

\sub Run the program using $\Delta t = .01$.  Now what is your
estimate for the duration of the epidemic?  What can you say about the
actual time required for $I$ to drop below 1?

\numa If we think the epidemic started with a single individual, we
can go backwards in time until $I$ is no longer greater than 1, and
see what the corresponding time is.  Do this for $\Delta t = -1, -.1,
\mbox{ and } -.01$.  What is your best estimate for the time that the
infection arrived?

\sub How many Recovereds were there at the start of the epidemic?
This would be the people who had presumably been infected in a
previous epidemic and now had immunity.

%%%%%%%%%%%%%%%%%%%%%%%%%%%%%%%%%%%%%%%%%%%%%%%%%%%%%%%%%%%%%%%%%%%%%%%
%                                                                     %
%                             Section 2                               %
%                                                                     % 
%%%%%%%%%%%%%%%%%%%%%%%%%%%%%%%%%%%%%%%%%%%%%%%%%%%%%%%%%%%%%%%%%%%%%%%

\section{The Mathematical Implications---\\Euler's Method}
\index{Euler's method}

\hyphenation{Eu-ler}
\subsection{Approximate Solutions}
\index{solution!approximate}

In the last section we approximated the function $S(t)$ by
piecewise-linear approximations using steps of size $\Delta t =1$, .1,
and .01. This process can clearly be extended to produce
approximations with an {\em arbitrary\/} number of steps. For any
given step size $\Delta t$, the result is a piecewise linear graph
whose segments are $\Delta t$ days wide. This graph then provides us
with an estimate for $S(t)$ for every value of $t$ in the interval $0
\leq t \leq 3$. We call this process of obtaining a function by
constructing a sequence of increasingly better approximations {\bf
  Euler's method}, after the Swiss mathematician Leonhard Euler
(1707--1783)\index{Euler, L.}. Euler was interested in the general problem
of finding the functions determined by a set of rate equations, and in
1768 he proposed this method to approximate them. The method is
conceptually simple and can indeed be used to get solutions for an
enormous range of rate equations.  For this reason we will make it a
basic tool.
    
To begin to get a sense of the general utility of Euler's method,
let's use it in a new setting.  Here is a simple problem that involves
just a single variable $y$ that depends on $t$.
\begin{align*}
\mbox{rate equation:}\quad y' &= 
	{\displaystyle .1\,y \left( 1-\frac{y}{1000} \right)}, \\
\mbox{initial condition:}\quad  y &= 100 \ 
		\mbox{when} \ t = 0.
\end{align*}
To solve the problem we must find the function $y(t)$ determined by
this rate equation and initial condition. We'll work without a
context, in order to emphasize the purely mathematical nature of
Euler's method. However, this rate equation is one member of a family
called {\bf logistic equations}\index{logistic
  equation}\label{logistic-begin} frequently used in population
models.  We will explore this context in the exercises.

Let's now construct the function that approximates $y(t)$ on the interval
$0 \leq t \leq 75$, using $\Delta t = .5$\,. We can make the suitable
modifications in SIRPLOT or SIRVALUE to get this
approximation. Suppose we want to view the approximation
graphically. Here's what the modified SIRPLOT would look like:

\begin{center}
{\bf Program:  modified SIRPLOT}
\end{center}

{\sl Set up GRAPHICS}
\vspace{-10pt}
\begin{verbatim}
   tinitial = 0
   tfinal = 75
   t = tinitial
   y = 100
   numberofsteps = 150
   deltat = (tfinal - tinitial)/numberofsteps
   FOR k = 1 TO numberofsteps
         yprime = .1 * y * (1 - y / 1000)
         deltay = yprime * deltat
\end{verbatim}
\vspace{-10pt}
\rule{50pt}{0pt} {\sl Plot the line from 
		{\tt (t, y)}} {\sl to 
		{\tt (t + deltat, y + deltay)}}\\
\vspace{-26pt}
\begin{verbatim}
         t = t + deltat
         y = y + deltay
   NEXT k
\end{verbatim}
As before, words in italics, like ``{\sl Plot the line from\ldots to}" need to
be translated into the specific formulation required by the computer
language you are using.

When we run this program, we get the following graph (axes and scales
have been added):

\vspace*{200pt}
\special{psfile=../ps/calc1/picture4.ps}\label{`logistic'}
  
\noindent
As before, what we see here is only an approximation
to the true solution $y(t)$.  How can we get some idea of how good
this approximation is?


\subsection{Exact Solutions}

For any chosen step size, we can produce an {\em approximate\/}
solution to a rate equation problem.  We will call such an
approximation an {\bf Euler approximation}. In the last section we saw
that we can improve the accuracy of the approximation by making the
steps smaller and using more of them.
	
For example, we have already found an approximate solution to the problem
\[
y' = {\displaystyle .1y \left( 1-\frac{y}{1000} \right)} \, ; 
    \qquad y(0) = 100
\]
on the interval $0 \leq t \leq 75$, using 150 steps of size $\Delta t
= 1/2$.  Consider a {\bf sequence}\index{sequence} of Euler
approximations to this problem that are obtained by increasing the
number of steps from one stage to the next.  To be systematic, let the
first approximation have 1 step, the next 2, the next 4, and so on.
(The important feature is that the number of steps increases from one
approximation to the next, not necessarily that they double---going up
by powers of 10 would be just as good.  A slight advantage in using
powers of 2 is to maximize computer accuracy.)  The number of steps
thus has the form $2^{j-1}$, where $j = $ 1, 2, 3, \ldots.  If we use
$y_j(t)$ to denote the approximating function with $2^{j-1}$ steps,
then we have an unending list:
\[
\begin{array}{cl}
	y_1(t)\, : & \mbox{Euler's approximation with 
		1 step} \\
	y_2(t)\, : & \mbox{Euler's approximation with 
		2 steps} \\
	y_3(t)\, : & \mbox{Euler's approximation with 
		4 steps} \\
	\vdots & \multicolumn{1}{c}{\vdots \rule{15pt}{0pt}} \\
	y_j(t)\, : & \mbox{Euler's approximation with 
		$2^{j - 1}$ steps} \\
	\vdots  &	\multicolumn{1}{c}{\vdots \rule{15pt}{0pt}}	
\end{array}
\]
Here are the graphs of $y_{j}(t)$ for $j=1,2,\ldots ,7$ plus the graph
of $y_{11}(t)$:

\vspace*{250pt}

\special{psfile=../ps/calc1/picture5.ps}\label{`sequence'}


These functions form a sequence of {\bf successive
  approximations}\index{successive approximation} to the true solution
$y(t)$, which is \index{solution!exact} obtained by taking the {\bf
  limit},\index{limit} as we did in the last section:
\[
y(t) = \lim_{j \to \infty} y_j(t).
\]
Earlier we noticed how the digits in the estimates for $S(3)$
stabilized.  If we
%
\mar{Functions and graphs can be limits, too}
%
plot the approximations $y_j(t)$ together we'll find that they
stabilize, too.  Each graph in the sequence is different from the
preceding one, but the differences diminish the larger $j$ becomes.
Eventually, when $j$ is large enough, the graph of $y_{j+1}$ does not
differ noticeably from the graph of $y_j$.  That is, the position of
the graph {\bf stabilizes}\index{stabilize} in the coordinate plane.
In this example, at the scale in the graph above, this happens around
$j=11$.  If we had drawn the graph of $y_{15}$ or $y_{20}$, it would
not have been distinguishable \mar{Euler's method is the process of
  finding solutions through a sequence of successive approximations}
from the graph of $y_{11}$.  It is this entire process of calculating
a sequence of successive approximations using increasingly many steps
as far as is needed to get the desired level of stabilization that is
meant when we talk about Euler's method.

The program SEQUENCE shown below plots 14 Euler approximations to
$y(t)$, increasing the number of steps by a factor of 2 each time.  It
demonstrates how the graphs of $y_j(t)$ stabilize to define $y(t)$ as
their limit.\index{program!SEQUENCE}


\begin{center}
{\bf Program: SEQUENCE \\ [2pt]
A sequence of graphs for \boldmath $y' = .1y(1-y/1000)$; $y(0) = 100$}
\end{center}

\noindent
{\sl Set up GRAPHICS}

\vspace{-10pt}
\begin{verbatim}
FOR j = 1 TO 14
\end{verbatim}
\vspace{-10pt}

\begin{minipage}[t]{3.8in}
\begin{verbatim}
tinitial = 0
tfinal = 75
t = tinitial
y = 100
numberofsteps = 2 ^ (j - 1)
deltat = (tfinal - tinitial) / numberofsteps
FOR k = 1 TO numberofsteps
      yprime = .1 * y * (1 - y / 1000)
      deltay = yprime * deltat
\end{verbatim}

\vspace{-12pt}
\rule{33pt}{0pt} {\sl Plot the line from 
		{\tt (t, y)}}\\[-1pt]
\rule{54pt}{0pt} {\sl to 
		{\tt (t + deltat, y + deltay)}}\\
\rule{33pt}{0pt} {\sl Color the line with color {\tt j}}

\vspace{-14pt}
\begin{verbatim}
      t = t + deltat
      y = y + deltay
NEXT k
\end{verbatim}
\end{minipage}
\begin{minipage}[t]{1.5in}
\mbox{}
\vspace{-6pt}

$
\left.
\rule{0pt}{108pt}
\right\} 
	\begin{array}{c} 
		\mbox{\bf Program:} \\ 
		\mbox{\bf modified SIRPLOT} 
	\end{array}
$

\end{minipage}

\vspace{-4pt}
\begin{verbatim}
NEXT j
\end{verbatim}

Notice that SEQUENCE contains the program SIRPLOT embedded in a loop
that executes SIRPLOT 14 times.  In this way SEQUENCE plots 14
different graphs.  The only new element that has been added to SIRPLOT
is ``{\sl Color the line with color {\tt j}}".  When you express this
in your programming language it instructs the computer to draw the
$j$-th graph using color number $j$ in the computer's ``palette.''  In
the exercises you are asked to use the program SEQUENCE to explore the
solutions to a number of rate equation problems.

%\medskip

\subsubsection{\bf Approximate solutions versus exact}  

By constructing successive approximations to the solution of a rate
equation problem, using a sequence of step sizes {\tt deltat} =
$\Delta t$ that shrink to 0, we obtain the {\em exact\/} solution in
the limit.

In practice, though, all we can ever get are particular
approximations. However, we can control the level of precision in our
approximations by adjusting the step size.  If we are dealing with a
model of some real process, then this is typically all we need.  For
example, when it comes to interpreting the $S$-$I$-$R$ model, we might
be satisfied to predict that there will be about 40500 susceptibles
remaining in the population after three days.  The table on
page~\pageref{backwards2} indicates we would get that level of
precision using a step size of about $\Delta t = 10^{-2}$.  Greater
precision than this may be pointless, because the modelling
process---which converts reality to mathematics---is itself only an
approximation.

The question we asked in the last section---In what sense do we know a
number?---applies equally to the functions we obtain using Euler's
method.  That is, even if we can characterize a function quite
precisely as the solution of a particular rate equation, we may be
able to evaluate it only approximately.
	
\subsection{A Caution}

We have now seen how to take a set of rate equations and find
approximations to the solution of these equations to any degree of
accuracy desired.  It is important to remember that all these
mathematical manipulations are only drawing inferences about the
model.  We are essentially saying that {\em if\/} the original
equations capture the internal dynamics of the situation being
modelled, {\em then\/} here is what we would expect to see.  It is
still essential at some point to go back to the reality being modeled
and check these predictions to see whether our original assumptions
were in fact reasonable, or need to be modified.  As Alfred North
Whitehead has said:
\begin{quote}
There is no more common error than to assume that, because prolonged
and accurate mathematical calculations have been made, the application
of the result to some fact of nature is absolutely certain.
\end{quote}
  
\newpage

\subsection{Exercises}

\subsubsection{Approximate solutions}

\vspace{-10pt}

\num  Modify SIRVALUE and SIRPLOT to analyze the population of 
Poland (see exercise~25 of chapter~1, page~\pageref{Poland}).  We assume the 
population $P(t)$ satisfies the conditions 
	$$
	P' = .009 \, P \quad \mbox{and} \quad 
		P(0) = \mbox{37,500,000},
	$$
where $t$ is years since 1985.  We want to know $P$ 100 
years into the future; you can assume that $P$ does not 
exceed 100,000,000.
\sub  Estimate the population in 2085.

\sub  Sketch the graph that describes this population growth.


\subsubsection{The Logistic Equation}
\index{logistic equation}

Suppose we were studying a population of rabbits.  If we turn 100
 rabbits loose in a field and let $y(t)$ be the number of rabbits at
 time $t$ measured in months, we would like to know how this function
 behaves. The next several exercises are designed to explore the
 behavior of the rate equation
\[
y' = {\displaystyle .1y \left( 1-\frac{y}{1000} \right)} \, ; 
    \qquad y(0) = 100
\]
and see why it might be a reasonable model for this system.

\num By modifying SIRVALUE in the way we modified SIRPLOT to get
 SEQUENCE, obtain a sequence of estimates for $y(37)$ that allows you
 to specify the exact value of $y(37)$ to two decimal places accuracy.

\numa Referring to the graph of $y(t)$ obtained in the text on
 page~\pageref{`sequence'}, what can you say about the behavior of $y$
 as $t$ gets large?

\sub Suppose we had started with $y(0)=1000$.  How would the
 population have changed over time? Why?

\sub Suppose we had started with $y(0)=1500$.  How would the
 population have changed over time? Why?

\sub Suppose we had started with $y(0)=0$.  How would the population
 have changed over time? Why?

\sub The number 1000 in the denominator of the rate equation is called
 the {\bf carrying capacity}\index{carrying capacity} of the system.
  Can you give a physical interpretation for this number?

\num Obtain graphical solutions for the rate equation for different
 values of the carrying capacity.  What seems to be happening as the
 carrying capacity is increased?  (Don't restrict yourself to $t = 37$
 here.)  In this problem and the next you should sketch the different
 solutions on the same set of axes.

\num Keep everything in the original problem unchanged except for the
 constant .1 out front.  Obtain graphical solutions with the value of
 this constant = .05, .2, .3, and .6.  How does the behavior of the
 solution change as this constant changes?

\num Returning to the original logistic equation, modify SIRVALUE or
 DO--WHILE to find the value for $t$ such that $y(t)=900$.

\num Suppose we wanted to fit a logistic rate equation to a
 population, starting with $y(0)=100$.  Suppose further that we were
 comfortable with the 1000 in the denominator of the equation, but
 weren't sure about the .1 out front. If we knew that $y(20)=900$,
 what should the value for this constant be?
\label{logistic-end} 

\subsubsection{Using SEQUENCE}
\index{program!SEQUENCE}

\vspace*{-10pt}

\num Each Euler approximation is made up of a certain number of
 straight line segments.  What instruction in the program SEQUENCE
 determines the number of segments in a particular approximation?  The
 first graph drawn has only a single segment.  How many does the fifth
 have? How many does the fourteenth have?

\num What is the slope of the first graph?  What are the slopes of the
 two parts of the second graph?  [You should be able to answer these
 questions without resorting to a computer.]

\num Modify the line in the program SEQUENCE which determines the
 number of steps by having it read {\tt numberofsteps = j}, and run
 the modified program.  Again, we are getting a sequence of
 approximations, with the number of steps increasing each time, but
 the approximations don't seem to be getting all that close to
 anything.  Explain why this modified program isn't as effective for
 our purposes as the original.

\num Modify SEQUENCE to produce a sequence of Euler approximations to
 the function $y(t)$ that satisfies the conditions
\[
y' = .2\, y(5 - y) \quad \mbox{and} \quad y(0) = 1.
\]
on the interval $0 \leq t \leq 10$.  [You need to change the final $t$
 value in the program, and you also need to ensure that the graphs
 will fit on your screen.]  \sub What is $y(10)$?  [If you add the
 line {\tt PRINT j, y} just before the line {\tt NEXT j}, a sequence
 of 14 estimates for $y(10)$ will appear on the screen with the
 graphs.]  \sub Make a rough sketch of the graph that is the limit of
 these approximations.  The right half of the limit graph has a
 distinctive feature; what is it?  \sub Without doing any
 calculations, can you estimate the value of $y(50)$?  How did you
 arrive at this value?  \sub Change the {\bf initial
   condition}\index{initial condition} from $y(0) = 1$ to $y(0) = 9$.
 Construct the sequence of Euler approximations beginning with {\tt
   numberofsteps} being 1, and make a rough sketch of the limit
 graph. What is $y(10)$ now?  Explain why the first several
 approximations look so strange.

\num Modify SEQUENCE to construct a sequence of Euler approximations
 for population of Poland (from exercise 1, above).  Sketch the limit
 graph $P(t)$, and mark the values of $P(0)$ and $P(100)$ at the two
 ends.

\num Construct a sequence of Euler approximations to the function
 $y(t)$ that satisfies the conditions
\[
y' = 2 \, t \quad \mbox{and} \quad y(0) = 0 
\]
over the interval $0 \leq t \leq 2$.  Note that this time the rate
 $y'$ is given in terms of $t$, not $y$.  Euler's method works equally
 well.  Using your sequence of approximations, estimate $y(2)$.  How
 accurate is your estimate?

\num Construct a sequence of Euler approximations to the function
 $y(t)$ that satisfies the conditions
\[
y' = \dfrac{4}{1+t^2} \quad \mbox{and} \quad y(0) = 0 
\]
over the interval $0 \leq t \leq 1$.  Estimate $y(1)$.  How 
accurate is your estimate?  [Note: the exact value of $y(1)$ 
is $\pi$, which your estimates may have led you to expect.  
By using special methods we shall develop much later we can 
prove that $y(1) = \pi$.]

%%%%%%%%%%%%%%%%%%%%%%%%%%%%%%%%%%%%%%%%%%%%%%%%%%%%%%%%%%%%%%%%%%%%%%%
%                                                                     %
%                             Section 3                               %
%                                                                     % 
%%%%%%%%%%%%%%%%%%%%%%%%%%%%%%%%%%%%%%%%%%%%%%%%%%%%%%%%%%%%%%%%%%%%%%%


\section{Approximate Solutions}
\index{solution!approximate}

Our efforts to find the functions that were determined by the rate
 equations for the $S$-$I$-$R$ model have brought to light several
 important issues:
\begin{itemize}

\item We often have to deal with a question that does not have a
 simple, straightforward answer; perhaps we are trying to determine a
 quantity (like the square root of~2, or $S(3)$ in the $S$-$I$-$R$
 model), to find some function (like $S(t)$), or to understand a
 process (like an epidemic, or buying and selling in a market).  An
 {\bf approximation} can get us started.

\item In many instances, we can make repeated improvements in the
 approximation.  If these {\bf successive approximations} get
 arbitrarily close to the unknown, and they do it quickly enough, that
 may answer the question for all practical purposes.  In many cases,
 there is no alternative.

\item The information that successive approximations give us is
 conveyed in the form of a {\bf limit}.

\item The method of successive approximations can be used to evaluate
 many kinds of mathematical objects, including numbers, graphs, and
 functions.

\item {\bf Limit processes} give us a valuable tool to probe difficult
 questions.  They lie at the heart of calculus.

\end{itemize}
     
\aside{Even the process of building a mathematical model for a
 physical system can be seen as an instance of successive
 approximations.  We typically start with a simple model (such as the
 $S$-$I$-$R$ model) and then add more and more features to it (e.g.,
 in the case of the $S$-$I$-$R$ model we might divide the population
 into different subgroups, have the parameters in the model depend on
 the season of the year, make immunity of limited duration, etc.).  Is
 it always possible, at least in theory, to get a sequence of
 approximating mathematical models that approaches reality in the
 limit?}


In the following chapters we will apply the process of successive
 approximation to many different kinds of problems.  For example, in
 chapter 3 the problem will be to get a better understanding of the
 notion of a rate of change of one quantity with respect to another.
  Then, in chapter 4, we will return to the task of solving rate
 equations using Euler's method.  Chapter 6 introduces the integral,
 defining it through a sequence of successive approximations.  As you
 study each chapter, pause to identify the places where the method of
 successive approximations is being used.  This can give you insight
 into the special role that calculus plays within the broader subject
 of mathematics.

To illustrate the general utility of the method, we end this chapter
 by returning to the problem raised in section~1 of constructing the
 values of $\sqrt{2}$ and $\pi$ to an arbitrary number of decimal
 places.

\subsection{Calculating $\pi$---The Length of a Curve}
\index{length of a curve}
\label{length-begin}

Humans were grappling with the problem of calculating $\pi$ at least
 3000 years ago.  In his work {\em Measurement of the Circle},
 Archimedes (287--212 {\sc b.c.}) used the method of successive
 approximations to calculate $\pi = 3.14 \ldots \, .$ He did this by
 starting with a circle of diameter 1, constructing an inscribed and a
 circumscribed hexagon, and calculating the lengths of their
 perimeters.  The perimeter of the circumscribed hexagon was clearly
 an overestimate for $\pi$, while the perimeter of the inscribed
 hexagon was an underestimate.  He then improved these estimates by
 going from hexagons to inscribed and circumscribed 12-sided polygons
 and again calculating the perimeters.  He repeated this process of
 doubling the number of sides until he had inscribed and circumscribed
 polygons with 96 sides. These left him with his final estimate
	$$
	3.1409\ldots = 
		3\frac{284\frac{1}{4}}{2017\frac{1}{4}} < \pi 
		< 3\frac{667\frac{1}{2}}{4673\frac{1}{2}} =
		3.1428\ldots
	$$  
	
In grade school we learned a nice simple formula for the length of a
 circle, but that was about it.  We were never taught formulas for the
 lengths of other simple curves like elliptic or parabolic arcs, for a
 very good reason---there are no such formulas.  There are various
 physical approaches we might take.  For example, we could get a rough
 approximation by laying a piece of string along the curve, then
 picking up the string and measuring it with a ruler.  Instead of a
 physical solution, we can use the essence of Archimedes' insight of
 approximating a circle by an inscribed ``polygon''---what we have
 earlier called a piecewise linear graph---to determine the length of
 any curve.  The basic idea is reminiscent of the way we made
 successive approximations to the functions $S(t)$, $I(t)$, and $R(t)$
 in the first section of this chapter.  Here is how we will approach
 the problem:
\begin{itemize}

\item approximate the curve by a chain of straight line segments;

\item measure the lengths of the segments;

\item use the sum of the lengths as an approximation to the true
 length of the curve.

\end{itemize}
Repeat this process over and over, each time using a chain that has
 shorter segments (and therefore more of them) than the last one.  The
 length of the curve emerges as the limit of the sums of the lengths
 of the successive chains.

\bigskip
\noindent
\fbox{
%\rule[-2.5in]{0in}{2.5in}\parbox{4.85in}{
\rule[-2.1in]{0in}{2.3in}
\parbox{4.85in}{
\begin{center}
\parbox{4.5in}{{\bf Distance Formula}\index{distance formula} If we
 are given two points $P_{1}(x_{1},y_{1})$ and $P_{2}(x_{2},y_{2})$ in
 the plane, then the distance between them is just
\begin{align*}
d&=\sqrt{{\Delta x}^2 + {\Delta y}^2} \\
  &=\sqrt{(x_{2}-x_{1})^2+(y_{2}-y_{1})^2}
\end{align*} 
That this follows directly from the Pythagorean
theorem\index{Pythagorean theorem}\label{Pythagorean theorem} can be
seen from the picture below:}
\end{center}
    
\setlength{\unitlength}{.57pt}
	\begin{picture}(210,90)(-180,75)
		\put(0,15){\vector(1,0){205}}
		\put(20,0){\vector(0,1){165}}
		\put(205,5){\makebox[0pt]{\small $x$}}
		\put(16,160){\makebox[0pt][r]{\small $y$}}
		\begin{thicklines}
			\put(10,50){\line(3,2){165}}
		\end{thicklines}
		\put(40,70){\circle*{4}}
		\put(160,150){\circle*{4}}
		\multiput(40,70)(2,0){60}{\circle*{1}}
		\multiput(160,70)(0,2){40}{\circle*{1}}
		\put(40,75){\makebox(0,0)[br]{\small $P_1$}}
		\put(160,155){\makebox(0,0)[br]{\small $P_2$}}
		\put(95,57){\makebox[0pt]{\small $\Delta x =x_{2}-x_{1}$}}
		\put(164,110){\makebox(0,0)[l]
			{\small $\Delta y=y_{2}-y_{1}$}}
	\end{picture}
}}


\vspace{.25in}


\smallskip
\noindent
\begin{minipage}[b]{2.6in}
\quad We'll demonstrate how this process works on a parabola.
  Specifically, consider the graph of $y = x^2$ on the interval $0
 \leq x \leq 1$.  At the right we have sketched the graph and our
 initial approximation.  It is a piecewise linear approximation with
 two segments whose end points have equally spaced $x$-coordinates.
\end{minipage}
\begin{picture}(130,115)(-35,0)
	\put(-10,0){\vector(1,0){130}}
	\put(0,-10){\vector(0,1){130}}
	\put(120,-10){\makebox[0pt]{$x$}}
	\put(-4,120){\makebox[0pt][r]{$y$}}
	\begin{thicklines}
		\bezier{200}(0,0)(50,0)(100,100)
		\put(0,0){\line(2,1){50}}
		\put(50,25){\line(2,3){50}}
	\end{thicklines}
	\put(104,97){\makebox[0pt][l]{$(1,1)$}}
	\put(56,19){\makebox[0pt][l]{$(.5, .25)$}}
	\put(83,57){\makebox[0pt][l]{$y = x^2$}}
	\put(0,0){\circle*{3}}
	\put(50,25){\circle*{3}}
	\put(100,100){\circle*{3}}
\end{picture}

\medskip

\enlargethispage*{10pt}
We can use the distance formula to find the lengths of the two segments.  
\begin{align*}
	\mbox{first segment}: & 
		\sqrt{(.5 - 0)^2 + (.25 - 0)^2} = .559016994
		 \\
	\mbox{second segment}: & 
		\sqrt{(1 - .5)^2 + (1 - .25)^2} = .901387819
\end{align*}
Their total length is the sum
	$$
	.559016994 + .901387819 = 1.460404813  .
	$$

The following program prints out the lengths of the two {\bf segments}
 and their {\bf total} length.\index{program!LENGTH}

\begin{center}
{\bf Program: LENGTH \\
Estimating the length of \boldmath $y = x^2$ over $0 \leq x \leq 1$}
\end{center}

\begin{verbatim}
   DEF fnf (x) = x ^ 2
   xinitial = 0
   xfinal = 1
   numberofsteps = 2
   deltax = (xfinal - xinitial) / numberofsteps
   total = 0
   FOR k = 1 TO numberofsteps
         xl = xinitial + (k - 1) * deltax
         xr = xinitial + k * deltax
         yl = fnf(xl)
         yr = fnf(xr)
         segment = SQR((xr - xl) ^ 2 + (yr - yl) ^ 2)
         total = total + segment
         PRINT k, segment
   NEXT k
   PRINT numberofsteps, total
\end{verbatim}
\label{length-end}


\subsection{Finding Roots with a Computer}
\index{root}

When we casually turn to our calculator 
%
\mar{Calculators and computers really work by making approximations}
%
and ask it for the value of $\sqrt{2}$, what does it really do?  Like
 us, the calculator can only add, subtract, multiply, and divide.
  Anything else we ask it to do must be reducible to these operations.
  In particular, the calculator doesn't really ``know'' the value of
 $\sqrt{2}$.  What it does know is how to approximate $\sqrt{2}$ to,
 say, 12 significant figures using only elementary arithmetic.  In
 this section we will look at two ways we might do this.  Apart from
 the fact that both approaches use successive approximations, they are
 remarkably different in flavor.  One works graphically, using a
 computer graphing package, and the other is a numerical algorithm
 that is about 4000 years old.

\subsubsection{A geometric approach}

Exercise 9 on page~\pageref{roots} considered the problem of finding
the roots of $f(x) = 1 - 2x^2$.  A bit of algebra confirms that
$\sqrt{2}/2$ is a root---i.e., $f(\sqrt{2}/2) = 0$.  The question is:
what is the {\em numerical value\/} of $\sqrt{2}/2$?

We'll answer this question by constructing a sequence of
 approximations that add digits, one at a time, to an estimate for
 $\sqrt{2}/2$.  Since the root lies at the point where the graph of
 $f$ crosses the $x$-axis, we just magnify the graph at this point
 over and over again, ``trapping'' the point between $x$ values that
 can be made arbitrarily close together.

\begin{picture}(320,160)(-5,-40)

%parabola graph

\put(0,0){\framebox(80,80){}}
\multiput(2,40)(2,0){39}{\circle*{1}}
\multiput(40,2)(0,2){39}{\circle*{1}}
\put(40,-12){\makebox[0pt]{$x$}}
\put(0,-12){\makebox[0pt]{\small -1}}
\put(80,-12){\makebox[0pt]{\small 1}}
\put(-4,37){\makebox[0pt][r]{$y$}}
\begin{thicklines}
	\bezier{200}(0,0)(40,160)(80,0)
\end{thicklines}
\put(64,36){\framebox(8,8){}}
\multiput(64,44)(3,3){12}{\circle*{1}}
\multiput(64,36)(3,-3){12}{\circle*{1}}

%first magnification

\put(140,90){\makebox[0pt]{\shortstack
	{first digit \\ [-2pt] determined}}}

\put(100,0){\framebox(80,80){}}
\multiput(102,40)(2,0){39}{\circle*{1}}
\put(140,-12){\makebox[0pt]{$x$}}
\put(100,-12){\makebox[0pt]{\small .60}}
\put(180,-12){\makebox[0pt]{\small .80}}
\begin{thicklines}
	\put(131.5,80){\line(1,-3){26.67}}
\end{thicklines}
\put(140,36){\framebox(8,8){}}
\put(139.5,36){\makebox(0,0)[tr]{\scriptsize .70}}
\put(148.2,36){\makebox(0,0)[tl]{\scriptsize .71}}
\multiput(140,44)(3,3){7}{\circle*{1}}
\multiput(140,36)(3,-3){7}{\circle*{1}}

%n-th magnification

\multiput(193,40)(7,0){3}{\circle*{2}}

\multiput(190,97)(7,0){3}{\circle*{2}}
\put(260,90){\makebox[0pt]{\shortstack
	{first {\em six\/} digits \\ [-2pt] determined}}}

\multiput(220,80)(-3,-3){5}{\circle*{1}}
\multiput(220,0)(-3,3){5}{\circle*{1}}
\put(220,0){\framebox(80,80){}}
\multiput(222,40)(2,0){39}{\circle*{1}}
\put(260,-12){\makebox[0pt]{$x$}}
\put(220,-12){\makebox[0pt]{\small .707100}}
\put(300,-12){\makebox[0pt]{\small .707110}}
\begin{thicklines}
	\put(263,80){\line(1,-3){26.67}}
\end{thicklines}
\put(272,36){\framebox(8,8){}}
\put(271.5,36){\makebox(0,0)[tr]{\scriptsize .707106}}
\put(309,36){\makebox(0,0)[tr]{\scriptsize .707107}}
\multiput(272,44)(3,3){7}{\circle*{1}}
\multiput(272,36)(3,-3){7}{\circle*{1}}

\multiput(313,40)(7,0){3}{\circle*{2}}

\multiput(310,97)(7,0){3}{\circle*{2}}

%caption

\put(150,-32){\makebox[0pt]{The graph of $y = 1 - 2x^2$ 
	under successive magnifications}}

\end{picture}


\noindent If we make each stage a ten-fold magnification over the
 previous one, then, as we zoom in on the next smaller interval that
 contains the root, one more digit in our estimate will be
 stabilized.\index{stabilize} The first six stages are described in
 the table below. They tell us $\sqrt{2}/2 = .707106\ldots$ to six
 decimal places accuracy.

\enlargethispage*{28pt}
Since this method of finding roots requires only that we be able to
 plot successive magnifications of the graph of $f$ on a computer
 screen, the method can be applied to any function that can be entered
 into a computer.
\begin{center}
The positive root of $1 - 2x^2$
\medskip

\begin{tabular}{cc}
	when the root lies between: & 
		\raisebox{-2pt}{the decimal expansion} \\
	\begin{tabular}{cc}
		lower value & upper value \\ \hline
		\begin{tabular}{l}
		.70 \rule{0pt}{14pt} \\
		.700 	\\
		.7070	\\
		.70710	\\
%		\rule{6pt}{0pt} \raisebox{2pt}{\vdots} \\
		.707100	\\
		.7071060 	\\
%		\rule{6pt}{0pt} \vdots 
		\rule{12pt}{0pt} \vdots 
		\end{tabular}
			&
		\begin{tabular}{l}
		.80 \rule{0pt}{14pt} \\
		.710 	\\
		.7080	\\
		.70720	\\
%		\rule{6pt}{0pt} \raisebox{2pt}{\vdots} \\
		.707110	\\
		.7071070 	\\
%		\rule{6pt}{0pt} \vdots 
		\rule{12pt}{0pt} \vdots 
		\end{tabular}
	\end{tabular}
		&
	\begin{tabular}{c}
		of the root begins with \\ \hline
		\begin{tabular}{l}
		.7 \rule{0pt}{14pt} \\
		.70 		\\
		.707		\\
		.7071	\\
%		\rule{6pt}{0pt} \raisebox{2pt}{\vdots} \\
		.70710	\\
		.707106 	\\
%		\rule{6pt}{0pt} \vdots 
		\rule{12pt}{0pt} \vdots 
		\end{tabular}
	\end{tabular}
\end{tabular}
\end{center}



\subsubsection{An algebraic approach -- the Babylonian algorithm}
\index{root}

About 4000 years ago Babylonian builders had a method for constructing
 the square root of a number from a sequence of successive
 approximations.  To demonstrate the method, we'll construct
 $\sqrt{5}$.  We want to find $x$ so that
\[
x^2 = 5 \qquad \mbox{or} \qquad x = \dfrac{5}{x}.
\]
The second expression may seem a peculiar way to characterize $x$, but
 it is at least equivalent to the first.  The advantage of the second
 expression is that it gives us {\em two\/} numbers to consider: $x$
 and 5/$x$.

For example, suppose we guess that $x = 2$.  Of course this is
 incorrect, because $2^2 = 4$, not 5.  The two numbers we get from the
 second expression are 2 and 5/2 = 2.5.  One is smaller than
 $\sqrt{5}$, the other is larger (because $2.5^2 = 6.25$).  Perhaps
 their {\em average\/} is a better estimate.  The average is 2.25, and
 $2.25^2 = 5.0625$.

Although 2.25 is not $\sqrt{5}$, it is a better estimate than either
 2.5 or 2.  If we change $x$ to 2.25, then
\[
x = 2.25 \, , \quad \dfrac{5}{x} = 2.222 \, , \quad
	\mbox{and their average is } \  2.236 \, .
\]
Is 2.236 a better estimate than 2.5 or 2.222?  Indeed it is: $2.236^2
 = 4.999696$.  If we now change $x$ to 2.236, a remarkable thing
 happens:
\[
x = 2.236 \, , \quad \dfrac{5}{x} = 2.236 \, , \quad
	\mbox{and their average is } \  2.236 \, !
\]
In other words, if we want accuracy to three decimal places, we have
 already found $\sqrt{5}$.  All the digits have stabilized.

Suppose we want {\em greater\/} accuracy?  Our routine readily
 obliges.  If we set $x$ = 2.236000, then
\[
x = 2.236000 \, , \quad \dfrac{5}{x} = 2.236136 \, ,
	 \quad \mbox{and their average is } \  2.236068\, .
\]
In fact, $\sqrt{5} = 2.236068\ldots$ is accurate to six decimal
 places.

Here is a summary of the argument we have just developed, expressed in terms of $\sqrt{a}$ for an arbitrary positive number $a$.  
\begin{center}
\begin{picture}(250,70)
	\begin{thicklines}
	\put(0,0){\framebox(250,60)
		{\shortstack{\bf If \boldmath $x$ is an estimate
			for \boldmath $\sqrt{a}$, \\
		\bf then the average of \boldmath $x$ and
			\boldmath $a/x$ \\ [2pt]
		\bf is a better estimate.}}}
	\end{thicklines}
\end{picture}
\end{center}

\noindent
Once we choose an {\em initial\/} estimate, this argument constructs a sequence of successive approximations to $\sqrt{a}$.  The process of constructing the sequence is called the {\bf Babylonian algorithm}\index{Babylonian algorithm}
 for square roots.  

A procedure that tells us how to carry out a sequence of steps, one at
a time, to reach a specific goal is called an {\bf
  algorithm}.\index{algorithm} Many algebraic processes are
algorithms. The word is a Latinization of the name of the Persian
astronomer Muhammad
al-Khw$\overline{\mbox{a}}$rizm$\overline{\mbox{\i}}$ (c.780\,{\sc
  a.d.}--c.850\,{\sc a.d.}), who lived and worked in Baghdad.  The title
of his seminal book {\em His}$\overline{\mbox{{\em a}}}${\em b al-jabr
  wal-muq}$\overline{\mbox{{\em a}}}$ {\em bala\/} (830\,{\sc
  a.d.})---usually referred to simply as {\em al-Jabr}---has a Latin
form that is even more familiar to mathematics students.

The Babylonian algorithm takes the current estimate $x$ for $\sqrt{a}$
 that we have at each stage and says ``replace $x$ by the average of
 $x$ and $a/x$.''  This kind of instruction is ideally suited to a
 computer, because {\tt A = B} in a computer program means ``replace
 the current value of {\tt A} by the current value of {\tt B}.''  In
 the program BABYLON printed below, the algorithm is realized by a
 {\tt FOR - NEXT} loop with a single line that reads ``{\tt x = (x + a
 / x) / 2}''.
	
The three-step procedure that we used in chapter 1 to obtain values of
 $S$, $I$, and $R$ in the epidemic problem is also an algorithm, and
 for that reason it was a straightforward matter to express it as the
 computer program SIRVALUE.
\begin{center}
{\bf Program: BABYLON\index{program!BABYLON} \\
\bf An algorithm to find \boldmath $\sqrt{a}$}
\end{center}

\begin{minipage}[t]{3.5in}
\begin{verbatim}
a = 5
x = 2
n = 6
FOR k = 1 TO n
      x = (x + a / x) / 2
      PRINT x
NEXT k
\end{verbatim}
\end{minipage}
\begin{minipage}[t]{1.2in}
Output:
\vspace{-6pt}
\begin{verbatim}
2.25
2.236111
2.236068
2.236068
2.236068
2.236068
\end{verbatim}
\end{minipage}

\subsection{Exercises}


In many of the questions below you are asked to do things like divide
 an arc into 2,000,000 segments or find a certain number to 8 decimal
 places.  This assumes you have fast computers \mar{Use common sense
 in deciding how close an approximation to make} available.  If you
 are using slower machines or programmable calculators, you should
 certainly feel free to scale back what is called for---perhaps to
 200,000 segments or 6 decimal places.  Use your own common sense;
 there's no value in sitting in front of a screen for an hour waiting
 for an answer to emerge.  To approximate a length by 200,000 segments
 will take 10 times as long as to approximate it by 20,000 segments,
 which in turn will take 10 times as long as the 2,000 segment
 approximation.  If you start with the cruder approximations, you
 should be able to get a good sense of what is reasonable to attempt
 with the facilities you have available, and modify what the problems
 call for accordingly.  \num By using a computer to graph $y = x^2 -
 2^x$, find the solutions of the equation $x^2 = 2^x$ to four decimal
 place accuracy.

\subsubsection{The program LENGTH}

\vspace{-10pt}

\num Run the program LENGTH to verify that it gives the lengths of the
 individual segments and their total length.

\num What line in the program gives the instruction to work with the
 function $f(x) = x^2$?  What line indicates the number of segments to
 be measured?

\num Each segment has a left and a right endpoint.  What lines in the
 program designate the $x$- and $y$-coordinates of the left endpoint;
 the right endpoint?

\num Where in the program is the length of the $k$-th segment
 calculated?  The segment is treated as the hypotenuse of a triangle
 whose length is measured by the Pythagorean theorem.  How is the base
 of that triangle denoted in the program?  How is the altitude of that
 triangle denoted?

\num Modify the program so that it uses 20 segments to estimate the
 length of the parabola.  What is the estimated value?

\num Modify the program so that it estimates the length of the
parabola using $200,\; 2,000,\; 20,000,\; 200,000,\; \mbox{ and }
2,000,000$ segments.  Compare your results with those tabulated below.
[To speed the process up, you will certainly want to delete the {\tt
    PRINT k, segment} statement that appears inside the loop.  Do you
  see why?]

\begin{center}
     \begin{tabular}{cc}\label{successive lengths}
		Number		  & Sum		\\
		of line segments & of their lengths \\[2pt]\hline
		\begin{tabular}{r}
   \rule{0pt}{15pt} 2    \\
                   20    \\
                  200    \\
               2\,000    \\
              20\,000    \\
             200\,000    \\
           2\,000\,000    
          \end{tabular}
		&        	
		\begin{tabular}{l}
		1.460\,404\,813  \rule{0pt}{15pt}  \\
		1.478\,756\,512    \\
		1.478\,940\,994    \\
		1.478\,942\,839    \\
		1.478\,942\,857    \\
		1.478\,942\,857    \\
		1.478\,942\,857    
		\end{tabular}
	\end{tabular}
\end{center}

\num What is the length of the parabola $y = x^2$ over the interval $0
 \leq x \leq 1$, correct to 8 decimal places?  What is the length,
 correct to 12 decimal places?

\num Starting at the origin, and moving along the parabola $y=x^2$,
 where are you when you've gone a total distance of 10?

\num Modify the program to find the length of the curve $y = x^3$ over
 the interval $0 \leq x \leq 1$.  Find a value that is correct to 8
 decimal places.

\num {\bf Back to the circle.}  Consider the unit circle centered at
 the origin.  Pythagoras' Theorem shows that a point $(x, y)$ is on
 the circle if and only if $x^2 + y^2 = 1$.  If we solve this for $y$
 in terms of $x$, we get $y=\pm \sqrt{1 - x^2}$, where the plus sign
 gives us the upper half of the circle and the minus sign gives the
 lower half. This suggests that we look at the function $g(x) =
 \sqrt{1 - x^2}$.  The arclength of $g(x)$ over the interval $-1 \leq
 x \leq 1$ should then be exactly $\pi$.

\sub Divide the interval into 100 pieces---what is the corresponding
 length?

\sub How many pieces do you have to divide the interval into to get an
 accuracy equal to that of Archimedes?

\sub Find the length of the curve $y = g(x)$ over the interval $-1
 \leq x \leq 1$, correct to eight decimal places accuracy.
 
\num This question concerns the function $h(x) = \dfrac{4}{1 + x^2}$.

\sub Sketch the graph of $y = h(x)$ over the interval $-2 \leq x \leq
 2$.

\sub Find the length of the curve $y = h(x)$ over the interval $-2
 \leq x \leq 2$.

\num Find the length of the curve $y=\sin x$ over the interval $0 \leq
 x \leq \pi$.


\subsubsection{The program BABYLON}

\vspace{-10pt}

\num Run the program BABYLON on a computer to verify the tabulated
 estimates for $\sqrt{5}$.

\num In whatever computer language you are using, it should be
 possible to tell the computer to print out more decimals.  Your
 teacher can tell you how this is handled on your computers.  Modify
 BABYLON to run with at least 14-digit precision in this and the
 following problems.  Also modify the program so it prints out the
 square of the estimate each time as well.  What is the estimated
 value of $\sqrt{5}$ in this circumstance?  How many steps were needed
 to get this value?  Use the square of this estimate as a measure of
 its accuracy.  What is the square?

\num Use the Babylonian algorithm to find $\sqrt{80}$.  

\sub First use 2 as your initial estimate.  How many steps are needed
 for the calculations to stabilize---that is, to reach a value that
 doesn't change from one step to the next?

\sub Since $9^2 = 81$, a good first estimate for $\sqrt{80}$ is 9.
  How many steps are needed this time for the calculations to
 stabilize?  Are the final values in (a) and (b) the same?

\num Use the Babylonian algorithm to find $\sqrt{250}$ and
 $\sqrt{1990}$.  If you use 2 as the initial estimate in each case,
 how many steps are needed for the calculations to stabilize?  If you
 use the integer nearest to the final answer as your initial estimate,
 then how many steps are needed?  Square your answers to measure their
 accuracy.

\num The Babylonian algorithm is considered to be {\bf very fast}, in
 the sense that each stage roughly doubles the number of digits that
 stabilize.  Does your work on the preceding exercises confirm this
 observation?  By comparison, is the routine that got the estimates
 for $S(3)$ (computed with the program SIRVALUE) faster or slower than
 the Babylonian algorithm?

%%%%%%%%%%%%%%%%%%%%%%%%%%%%%%%%%%%%%%%%%%%%%%%%%%%%%%%%%%%%%%%%%%%%%%%
%                                                                     %
%                             Section 4                               %
%                                                                     % 
%%%%%%%%%%%%%%%%%%%%%%%%%%%%%%%%%%%%%%%%%%%%%%%%%%%%%%%%%%%%%%%%%%%%%%%

\section{Chapter Summary}

\subsection{The Main Ideas}

\begin{itemize}

\item The exact numerical value of a quantity may not be known; the
 value is often given by an {\bf approximation}.

\item A numerical quantity is often given as the {\bf limit} of a
 sequence of {\bf successive approximations}.

\item When a particular digit in a sequence of successive
 approximations {\bf stabilizes}, that digit is assumed to appear in
 the limit.

\item {\bf Euler's method} is a procedure to construct a sequence of
 increasingly better {\bf approximations} of a function defined by a
 set of rate equations and initial conditions.  Each approximation is
 a piecewise linear function.

\item The exact function defined by a set of rate equations and
 initial conditions can be expressed as the {\bf limit} of a sequence
 of {\bf successive Euler approximations} with smaller and smaller
 step sizes.

\end{itemize}

\subsection{Expectations}

\begin{itemize}

\item You should be able to use a program to {\bf construct a
 sequence} of estimates for $S$, $I$, and $R$, given a specific
 $S$-$I$-$R$ model with initial conditions.

\item You should be able to {\bf modify} the SIRVALUE and SIRPLOT
 programs to construct a sequence of estimates for the values of
 functions defined by other rate equations and initial conditions.

\item You should be able to use programs that construct a sequence of
 {\bf Euler approximations} for the function defined by a rate
 equation with an initial condition.  The programs should provide both
 tabular and graphical output.

\item You should be able to estimate the values of the roots of an
 equation $f(x) = 0$ using a {\bf computer graphing package}.

\item You should be able to find a square root using the {\bf
 Babylonian algorithm}.

\item You should be able to find the length of any piece of any curve.

\end{itemize}

\subsection{Chapter Exercises}

\numa  We have considered the logistic equation 
\[
y' = {\displaystyle .1y \left( 1-\frac{y}{1000} \right)}
\]
On page~\pageref{`logistic'} we looked at the resulting graph of $y$
 vs. $t$, using the starting value $y(0)=100$.  There is another
 graphical interpretation, though, that is also instructive.  Note
 that the logistic equation specifies $y'$ as a function of~$y$.
  Sketch the graph of this function.  That is, plot $y$ values on the
 horizontal axis and $y'$ values on the vertical axis; points on the
 graph will thus be of the form $(y, y')$, where the $y'$-coordinate
 is the value given by the logistic equation for the given $y$ value.
  What shape does the graph have?

\sub For what value of $y$ does this graph take on its largest value?
  Where does this $y$ value appear in the graph of $y(t)$ versus $t$?

\sub For what values of $y$ does the graph of $y'$ versus $y$ cross
 the $y$-axis?  Where do these $y$ values appear in the graph of
 $y(t)$ versus $t$?

\subsubsection{Grids on Graphs}

When writing a graphing program it is often useful to have the
 computer draw a grid on the screen.  This makes it easier to estimate
 numerical values, for instance.  We can use a simple FOR--NEXT loop
 inside a program to do this.  For instance, suppose we had written a
 program (SIRPLOT or SEQUENCE, for instance) with the graphics already
 in place.  Suppose the screen window covered values from 0 to 100
 horizontally, and 0 to 50,000 vertically.  If we wanted to draw 21
 vertical lines (including both ends) spaced 5 units apart and 11
 horizontal lines spaced 5,000 units apart, the following two loops
 inserted in the program would work:

\newpage

%\noindent
\begin{verbatim}
FOR k = 0 TO 20
\end{verbatim}
\vspace{-12pt}
\rule{34pt}{0pt} {\sl Plot the line from 
		{\tt (5 * k, 0)}} {\sl to {\tt (5 * k, 50000)}}\\
\vspace{-25pt}
\begin{verbatim}
NEXT k
FOR k = 0 TO 10
\end{verbatim}
\vspace{-12pt}
\rule{34pt}{0pt} {\sl Plot the line from 
		{\tt (0, 5000 * k)}} {\sl to {\tt (100, 5000 * k)}}\\
\vspace{-25pt}
\begin{verbatim}
NEXT k
\end{verbatim}
You should make sure you see how these loops work and that you can
 modify them as needed.

\num If the screen window runs from $-20$ to 120 horizontally, and 250
 to 750 vertically, how would you modify the loops above to create a
 vertical grid spaced 10 units apart and a horizontal grid spaced 25
 units apart?

\num Go back to our basic $S$-$I$-$R$ model.  Modify SIRPLOT to
 calculate and plot on the same graph the values of $S(t)$, $I(t)$,
 and $R(t)$ for $t$ going from 0 to 120, using a stepsize of $\Delta t
 = .1$.  Include a grid with a horizontal spacing of 5 days, and a
 vertical spacing of 2000 people.  You should get something that looks
 like this:


\vspace{240pt}

\special{psfile =../ps/calc1/sir.ps}
 
 
	


